\begin{abstract}
  This study investigates the effects of \textit{fine-tuning} on language models tasked with navigating ASCII mazes. By analyzing activation graphs generated from model embeddings, we assess structural changes using various graph metrics with distribution inference through statistical sampling. Our findings indicate subtle yet non-significant alterations in graph structures post-\textit{fine-tuning}, suggesting that while \textit{fine-tuning} influences information processing, these changes are not drastically reflected in the analyzed metrics. The study highlights the importance of considering different abstraction levels within models to fully understand the impact of \textit{fine-tuning}. Future recommended works include exploring diverse prompts, activation evolution during task resolution, and the effects of various \textit{fine-tuning} techniques and model architectures on activation graph structures.
\end{abstract}

\begin{resumo} 
  Este estudo investiga os efeitos do \textit{fine-tuning} em modelos de linguagem encarregados de navegar por labirintos em ASCII. Ao analisar os grafos de ativações gerados a partir dos embeddings dos modelos, avaliamos as mudanças estruturais utilizando diversas métricas de grafos com inferência de distribuições por amostragem estatística. Nossos achados indicam alterações sutis, porém não significativas, nas estruturas dos grafos após o \textit{fine-tuning}, sugerindo que, embora o \textit{fine-tuning} influencie o processamento de informações, essas mudanças não são drasticamente refletidas nas métricas analisadas. O estudo destaca a importância de considerar diferentes níveis de abstração dentro dos modelos para compreender plenamente o impacto do \textit{fine-tuning}. Trabalhos futuros recomendados incluem a exploração de prompts diversos, a evolução das ativações durante a resolução de tarefas e os efeitos de várias técnicas de \textit{fine-tuning} e arquiteturas de modelos nas estruturas dos grafos de ativações.  
\end{resumo}